\section*{Capítulo \ref{ch:g11:drift-to-spreading} --- Da deriva continental à expansão do assoalho oceânico}

\begin{solution}{exo:g11:drift-to-spreading:1}
A \omterm{def:g11:drift-to-spreading:spreading}{expansão do assoalho oceânico} é a fabricação de novo assoalho oceânico na
crista de uma dorsal oceânica, onde a rocha fundida sobe e se solidifica,
seguida do seu movimento para longe da dorsal dos dois lados. A taxa total é
a velocidade com que as duas placas se afastam uma da outra; a semitaxa é a
velocidade com que o assoalho de um lado se afasta do eixo. Para uma
expansão simétrica, taxa total $= 2
\times$ semitaxa.
\end{solution}

\begin{solution}{exo:g11:drift-to-spreading:2}
Acima da sua \omterm{def:g11:drift-to-spreading:remanence}{temperatura de Curie}, cerca de \qty{580}{\celsius}, a magnetita
não consegue manter uma magnetização própria; só quando esfria abaixo dela os
seus grãos se magnetizam ao longo do campo e conservam essa direção.
Aquecida de novo a \qty{700}{\celsius}, acima da \omterm{def:g11:drift-to-spreading:remanence}{temperatura de Curie}, a
rocha perde o seu registro antigo; ao esfriar de novo, ela registra o campo
dessa nova época.
\end{solution}

\begin{solution}{exo:g11:drift-to-spreading:3}
O som desce e volta: $d = 1500 \times 5.2/2 = \qty{3900}{m}$.
\end{solution}

\begin{solution}{exo:g11:drift-to-spreading:4}
O intervalo normal atual começou há \num{0.78} milhão de anos: ele já dura
\num{780000} anos. O evento Jaramillo durou $1.07 - 0.99 = \num{0.08}$
milhão de anos, \num{80000} anos. Contando as fronteiras com menos de cinco
milhões de anos (0.78, 0.99, 1.07, 1.77, 1.95, 2.14, 2.15, 2.581, 3.04, 3.11,
3.22, 3.33, 3.58, 4.18, 4.29, 4.48, 4.62, 4.80, 4.89, 4.98): vinte inversões.
\end{solution}

\begin{solution}{exo:g11:drift-to-spreading:5}
$65\,000/40\,000 \approx 1.6$: as dorsais dariam cerca de uma volta e meia
à Terra.
\end{solution}

\begin{solution}{exo:g11:drift-to-spreading:6}
$20/0.78 \approx \qty{25.6}{km}$ por milhão de anos, isto é, uma semitaxa de
cerca de \qty{2.6}{cm} por ano; a taxa total é de cerca de \qty{5.1}{cm} por
ano.
\end{solution}

\begin{solution}{exo:g11:drift-to-spreading:7}
Oeste: $45/3.58 \approx 12.6$~km por milhão de anos, \qty{1.26}{cm} por ano;
leste: $44/3.58 \approx 12.3$~km por milhão de anos, \qty{1.23}{cm} por ano.
Taxa total: $1.26 + 1.23 \approx \qty{2.5}{cm}$ por ano. As duas semitaxas
diferem cerca de \qty{2}{\percent}: a expansão é quase exatamente simétrica.
\end{solution}

\begin{solution}{exo:g11:drift-to-spreading:8}
\qty{1.9}{cm} por ano são \qty{19}{km} por milhão de anos: $990/19 \approx 52$
milhões de anos, um pouco mais que os 49 milhões de anos do sedimento mais
antigo (cerca de \qty{6}{\percent}).
\end{solution}

\begin{solution}{exo:g11:drift-to-spreading:9}
A \qty{20}{km} por milhão de anos: \qty{15.6}{km}, \qty{19.8}{km},
\qty{21.4}{km}, \qty{35.4}{km} e \qty{39.0}{km}. A partir do eixo: normal de
0 a \qty{15.6}{km}, inversa até \qty{19.8}{km}, normal até \qty{21.4}{km}
(Jaramillo), inversa até \qty{35.4}{km}, normal até \qty{39.0}{km}
(Olduvai), depois inversa.
\end{solution}

\begin{solution}{exo:g11:drift-to-spreading:10}
$64 \times 10^6/4.6 \times 10^9 \approx \num{0.014}$, cerca de
\qty{1.4}{\percent} da idade da Terra. Isso apoia o argumento de Hess de que
o assoalho oceânico é jovem porque está sempre sendo fabricado e destruído:
nenhum assoalho oceânico sobreviveu da maior parte da história da Terra.
\end{solution}

\begin{solution}{exo:g11:drift-to-spreading:11}
A cada momento, o assoalho se forma no eixo com a polaridade desse momento,
e se divide: uma metade se afasta de cada lado com a mesma velocidade. Duas
faixas formadas ao mesmo tempo ficam, portanto, à mesma distância do eixo,
uma de cada lado, com a mesma polaridade: cada lado é uma cópia do outro,
refletida no eixo. A simetria é quebrada se o assoalho se afasta mais
depressa de um lado que do outro (expansão assimétrica), ou se o eixo salta
para uma nova posição.
\end{solution}

\begin{solution}{exo:g11:drift-to-spreading:12}
(1) \qty{5}{cm} são $5 \times 10^{-5}$~\unit{km}: a cada ano, acrescenta-se
uma faixa de $65\,000 \times 5 \times 10^{-5} = \qty{3.25}{km^2}$.
(2) $2.9 \times 10^8/3.25 \approx 9 \times 10^7$ anos, cerca de 90 milhões
de anos. (3) Isso é da mesma ordem de grandeza que a idade do assoalho mais
antigo (190 milhões de anos), e menor que ela: para que o assoalho oceânico
continue tão jovem, é preciso que tanto assoalho seja destruído, em algum
lugar, quanto é fabricado nas dorsais.
\end{solution}

\begin{solution}{exo:g11:drift-to-spreading:13}
O evento durou \num{0.08} milhão de anos. A \qty{12.5}{km} por milhão de
anos, a sua listra tem $0.08 \times 12.5 = \qty{1}{km}$ de largura; a
\qty{75}{km} por milhão de anos, $0.08 \times 75 = \qty{6}{km}$. Vista da
superfície do mar, \qty{3}{km} acima, uma listra muito mais estreita que a
sua profundidade só produz um campo fraco, confundido com os campos das suas
largas vizinhas de polaridade oposta; uma listra duas vezes mais larga que a
sua profundidade se destaca.
\end{solution}

\begin{solution}{exo:g11:drift-to-spreading:14}
Local 16: $176/11 \approx \qty{16}{m}$ por milhão de anos; local 20: $72/67
\approx \qty{1.1}{m}$ por milhão de anos. O sedimento se acumula a taxas
muito diferentes de um lugar para outro e de uma época para outra, de modo
que a sua espessura não dá a idade do assoalho. A idade é dada pelos
fósseis do sedimento mais antigo, assentado diretamente sobre o basalto.
\end{solution}

\begin{solution}{exo:g11:drift-to-spreading:15}
(1) Oeste: $58/2.581 \approx \qty{22.5}{km}$ por milhão de anos,
\qty{2.25}{cm} por ano; leste: $70/2.581 \approx \qty{27.1}{km}$ por milhão
de anos, \qty{2.71}{cm} por ano; taxa total $128/2.581 \approx \qty{49.6}{km}$ por
milhão de anos, cerca de \qty{5.0}{cm} por ano. (2) Cerca de \qty{225}{km} a
oeste e \qty{271}{km} a leste. (3) As duas fronteiras estavam juntas no eixo
quando a inversão aconteceu; desde então, elas se afastaram à taxa total,
qualquer que seja a divisão do movimento entre os dois lados. A sua
separação dividida pela sua idade é a taxa total.
\end{solution}

\begin{solution}{pb:g11:drift-to-spreading:1}
\textbf{1.} O local 16 é o mais próximo (\qty{191}{km}, 11 milhões de
anos), o local 20 o mais distante (\qty{1270}{km}, 67 milhões de anos):
quanto mais longe da dorsal, mais antigo o assoalho.

\textbf{2.} O sedimento começa a se depositar sobre o basalto assim que o
basalto se forma, de modo que a camada mais antiga tem mais ou menos a idade
do basalto. Ela só pode ser um pouco mais jovem, nunca mais velha.

\textbf{3.} Em quilômetros por milhão de anos: local 16:
$191/11 \approx 17.4$; local 15: $380/24 \approx 15.8$; local 18: $506/26 \approx 19.5$;
local 17: $643/33 \approx 19.5$; local 14: $727/40 \approx 18.2$; local 19: $990/49
\approx 20.2$; local 20: $1270/67 \approx 19.0$.

\textbf{4.} $129.5/7 \approx \num{18.5}$~km por milhão de anos, cerca de
\qty{1.85}{cm} por ano.

\textbf{5.} Sim: as suas razões (19.5) ficam entre as dos locais do oeste. O
assoalho se expandiu mais ou menos à mesma taxa dos dois lados: a expansão é
simétrica.

\textbf{6.} $t = d/u$: a idade é proporcional à distância.

\textbf{7.} $u = 208\,276/10\,952 \approx \num{19.0}$~km por milhão de anos,
isto é, \qty{1.90}{cm} por ano.

\textbf{8.} Cerca de $2 \times 1.9 = \qty{3.8}{cm}$ por ano.

\textbf{9.} Sim: \qty{1.9}{cm} por ano é um pouco menos que \qty{2}{cm} por
ano.

\textbf{10.} $1270/19.0 \approx 67$ milhões de anos, o mesmo que o seu
sedimento mais antigo.

\textbf{11.} $191/9 \approx \num{21.2}$~km por milhão de anos,
\qty{2.1}{cm} por ano.

\textbf{12.} $1270/71 \approx \num{17.9}$~km por milhão de anos,
\qty{1.8}{cm} por ano.

\textbf{13.} $180\,329/9854 \approx \num{18.3}$~km por milhão de anos,
\qty{1.83}{cm} por ano.

\textbf{14.} $(1.90 - 1.83)/1.83 \approx \qty{4}{\percent}$: as duas taxas
concordam dentro das incertezas das idades e das distâncias.

\textbf{15.} Uma idade vem dos fósseis do sedimento, a outra do magnetismo
do basalto e da escala de tempo das inversões. Um erro num método não
apareceria no outro: a sua concordância é um teste forte da \omterm{def:g11:drift-to-spreading:spreading}{expansão do
assoalho oceânico}.

\textbf{16.} $13.8 + 13.8 = \ang{27.6}$.

\textbf{17.} $111.2 \times \cos 30^\circ \approx \qty{96.3}{km}$ por grau;
$27.6 \times 96.3 \approx \qty{2660}{km}$.

\textbf{18.} $2660/19.0 \approx 140$ milhões de anos.

\textbf{19.} $2660/130 \approx \num{20.5}$~km por milhão de anos, cerca de
\qty{2.0}{cm} por ano, um pouco mais que \qty{1.9}{cm} por ano. Como os
últimos 67 milhões de anos correram a cerca de \qty{1.9}{cm} por ano, a
abertura inicial foi um pouco mais rápida: o assoalho entre o local 20 e a
borda, $2660 - 1270 =
\qty{1390}{km}$, formou-se em $130 - 67 = 63$ milhões de anos, cerca de
\qty{2.2}{cm} por ano.

\textbf{20.} \textbf{O assoalho oceânico mais antigo junto à África a
\ang{30} Sul tem cerca de 130 milhões de anos} (140 milhões de anos
extrapolando a taxa atual): o Atlântico Sul começou a se abrir ali no
período Cretáceo.
\end{solution}
